%\section{Basis calculus for substitution structure}

\section{Substitution Structure}

We here investigate the structure of worlds in \firstOrderLogic{}s and how they can be represented by tensors.
As we already argued, groundings of sentences can be aligned in boolean tensors, to be called grounding tensors.
We here investigate sparse representation schemes for this \substitutionStructure{}.
In particular we apply basis calculus to find sparse decompositions of the grounding tensors


\red{The semantic structure, an attempt to represent in tensors multiple worlds at once, is to be investigated in \charef{cha:probabilisticRelationalModels}.}


%% From Report
%We in this section investigate the structure of terms and formulas in a single \firstOrderLogic{} world, which we for now index by $\worldindices$.
%As we have derived in \charef{cha:logicalRepresentation} for \propositionalLogic{}, we develop efficient tensor network representations of the representing tensors based on the corresponding syntax.

\subsection{Grounding Tensors}

Given a \firstOrderLogic{} world $\worldindices$, arbitrary formulas are interpreted in terms of the satisfactions of their groundings.
We define their semantic first, and then relate their syntactical decomposition to tensor networks, similar to our approach to \propositionalLogic{} in \charef{cha:logicalRepresentation}.

\begin{definition}[Grounding of a first-order formula given a world]
    Given a specific world $\worldindices$, with a domain $\worlddomain$ enumerated by $[\inddim]$, the grounding of a formula $\folexformula$ with variables $\indvariableof{\folexformula}$  is the tensor
    \begin{align*}
        \groundingofat{\folexformula}{\indvariableof{\folexformula}} :
        \bigtimes_{\indenumerator\in[\indvariableof{\folexformula}]} [\inddim] \rightarrow \ozset \, .
    \end{align*}
    Each coordinate represents thereby the boolean, whether the substitution of the variables in the formula is satisfied given a world $\worldindices$, that is
    \begin{align*}
        \groundingofat{\folexformula}{\indexedindvariableof{\folexformula}} = 1
    \end{align*}
    if and only if the substitution of $\folexformula$ with the variables $\indvariableof{\folexformula}$ replaced by the objects $\indexinterpretationat{\indindexof{\indenumerator}}$ is satisfied on the world $\worldindices$.
\end{definition}

\todo{We in the following investigate how these tensors are represented by tensor networks involving the world defining tensors.}


%% Basis encoding
%When interpreting this map as a basis encoding, formulas are tensors in the tensor space
%\begin{align*}
% 	\left(\bigotimes_{\atomenumeratorin, \indindexlist\in[\inddim]} \rr^{2} \right) \otimes
%	\left(  \bigotimes_{\atomenumeratorin, \indindexof{0},\ldots,\indindexof{\indorder_{\folexformula}}\in[\inddim]} \rr^{2} \right) \, .
%\end{align*}

%\subsection{Terms}
%
%Terms are object valued maps on $\worlddomain^{n}$.
%The basis encoding of each term is a boolean and directed tensor.
%% Constants: Functions with $n=0$
%Constants are the functions with $n=0$, and are represented by basis vectors in $\rr^{\cardof{\worlddomain}}$.
%Given database semantics, this basis vector does not vary among different worlds.
%% General terms
%Most general terms are combinations of functions and constants and their basis encodings are acyclic contractions of basis encodings to functions.

\subsection{Basis encodings of terms}

Terms are object valued functions, which $r_{\folterm}$ different arguments are also object valued and called variables.
We represent terms by the corresponding basis encodings
\begin{align*}
    \bencodingofat{\folterm}{
        \indvariableof{\folterm}, \indvariableof{[r_{\folterm}]}
    }
\end{align*}

Basic blocks:
\begin{itemize}
    \item Constant
    \stantikzpic{
        \draw[->-] (0,1) -- (0,2.5) node[midway,right] {\colorlabelsize $\indvariableof{\folterm}$};
        \drawsquaretensor{0}{0}{$\onehotmapof{\indindex}$}
    }
    \item Variable
    \stantikzpic{
        \draw[->-] (0,1) -- (0,2.5) node[midway,right] {\colorlabelsize $\indvariableof{\folterm}$};
        \drawsquaretensor{0}{0}{$\identity$}
        \draw[->-] (0,-2.5) -- (0,-1) node[midway,right] {\colorlabelsize $\indvariable$};
    }
    \item Function of terms
    \stantikzpic{
        \draw[->-] (2,1) -- (2,2.5) node[midway,right] {\colorlabelsize $\indvariableof{\folterm}$};
        \draw (-1,-1) rectangle (5,1);
        \drawtextnode{2}{0}{\corelabelsize $\bencodingof{\exfunction}$}
        \draw[->-] (0,-2.5) -- (0,-1) node[midway,right] {\colorlabelsize $\indvariableof{0}$};
        \draw[->-] (4,-2.5) -- (4,-1) node[midway,right] {\colorlabelsize $\indvariableof{r_{\exfunction}}$};
    }
\end{itemize}
Contractions of those terms along directed acyclic graphs with single root nodes represent more complex terms.
As such, basis encodings of functions can be contracted with constants, variables or other functions to build up basis encodings of more complex terms.

Example: Father of the father of alice is represented by the contraction
\stantikzpic{
    \draw[->-] (2,1) -- (2,2.5) node[midway,right] {\colorlabelsize $\indvariableof{\folterm}$};
    \draw (-1,-1) rectangle (5,1);
    \drawtextnode{2}{0}{\corelabelsize $\bencodingof{Father}$}
    \draw[->-] (2,-2.5) -- (2,-1) node[midway,right] {\colorlabelsize $\indvariable$};
    \draw (-1,-4.5) rectangle (5,-2.5);
    \drawtextnode{2}{-3.5}{\corelabelsize $\bencodingof{Father}$}
    \draw[->-] (2,-6) -- (2,-4.5) node[midway,right] {\colorlabelsize $\indvariable$};

    \drawsquaretensor{2}{-7}{$\onehotmapof{Alice}$}
}


% Contraction of world defining 
We notice, that the basis encoding of each term is decomposed into a network of world defining tensors, namely constant and function defining tensors.

\subsection{Atomic sentences}

Atomic sentences are contractions of coordinate encodings of predicates with basis encodings of terms.

Example: Knowing the father of alice
\stantikzpic{
%    \draw[->-] (2,1) -- (2,2.5) node[midway,right] {\colorlabelsize $\indvariableof{\folexformula}$};
    \draw (-1,-1) rectangle (5,1);
    \drawtextnode{2}{0}{\corelabelsize ${Knows}$}
    \draw[->-] (0,-2.5) -- (0,-1) node[midway,right] {\colorlabelsize $\indvariableof{0}$};
    \draw[->-] (4,-2.5) -- (4,-1) node[midway,right] {\colorlabelsize $\indvariableof{1}$};

    \draw (2,-4.5) rectangle (6,-2.5);
    \drawtextnode{4}{-3.5}{\corelabelsize $\bencodingof{Father}$}
    \draw[->-] (4,-6) -- (4,-4.5) node[midway,right] {\colorlabelsize $\indvariable$};

    \drawsquaretensor{4}{-7}{$\onehotmapof{Alice}$}
}


%\subsection{Complex sentences}
\todo{Complex sentences are built by
\begin{itemize}
    \item Quantifier free sentences are contractions of basis encodings to atomic sentences with basis encodings of the connectives.
    \item Quantifiers act as non-linear transforms of tensors.
\end{itemize}
}


%\subsection{Substitution by slicing}
% Slicing interpretation
%Slicing the grounding tensor of a formula a first-order formula amounts to substitution of the respective variable by the constant at the enumeration index.

%\subsection{Syntactical Decomposition of quantifier-free formulas}

\subsection{Formula Synthesis by Connectives}\label{sec:folConnectiveRepresentation}

In order to have a sound semantic, the grounding of \firstOrderLogic{} formulas is determined by the syntax of the formula, i.e. a decomposition of the formula into connectives and quantifiers acting on atomic formulas.

% Formulas as maps from worlds to groundings
Quantifier-free formulas are connectives acting on atomic formulas.
We can describe them as in the case of \propositionalLogic{} in the $\bencodingof{}$-formalism.
While the atomic formulas where delta tensors copying states, they are more involved here.

\begin{theorem}
    For any connective $\exconnective$ and formulas $\folexformula_1$ and $\folexformula_2$ we have
    \begin{align}
        &\groundingofat{(\folexformula_1\exconnective\folexformula_2)}{\indvariableof{\folexformula_1}\cup\indvariableof{\folexformula_2}} \\
        &\quad=
        \contractionof{
            \bencodingofat{\groundingof{\folexformula_1}}{\headvariableof{\folexformula_1},\indvariableof{\folexformula_1}},
            \bencodingofat{\groundingof{\folexformula_2}}{\headvariableof{\folexformula_2},\indvariableof{\folexformula_2}},
            \bencodingofat{\exconnective}{\headvariableof{\folexformula_1\exconnective\folexformula_2}, \headvariableof{\folexformula_1}, \headvariableof{\folexformula_2}},
            \tbasisat{\headvariableof{\folexformula_1\exconnective\folexformula_2}}
        }
        {\shortindvariablelist} \, .
    \end{align}
\end{theorem}
\begin{proof}
    This directly follows from \theref{the:compositionByContraction}.
%	By the semantic interpretation of the groundings, which has to be sound.
\end{proof}

% Shared variables
Here, variables can be shared by the connected formulas, therefore the variables in the combined formula are unions of the possible not disjoint variables of the connected formulas.

%% Propositional interpretation
%When we understand the head variables in the basis encoding of atoms as the categorical variables, and get a similar interpretation of the tensor network decomposition as in the propositional case.
%\subsection{Propositionalization}

When interpreting the head variables of relational encoded atomic formulas as the atoms of a propositional theory, we find a propositional formula $\exformula$ associated with any decomposable \firstOrderLogic{} formula.

\begin{definition}
    \label{def:propositionalEquivalent}
    Given a formula $\folexformula$ in \firstOrderLogic{}, we say that a propositional formula $\formulaat{\shortcatvariables}$ is the propositional equivalent to $\folexformula$ given atomic formulas $\extformulaof{\atomenumerator}$ in \firstOrderLogic{}, when for any world $\worldindices$ we have
    \begin{align*}
        \groundingofat{\folexformula}{\indvariableof{\folexformula}}
        = \contractionof{
            \{\bencodingofat{\groundingof{\extformulaof{\atomenumerator}}}{\catvariableof{\atomenumerator},\indvariableof{\extformulaof{\atomenumerator}}} : \atomenumeratorin\}
            \cup \{\formulaat{\shortcatvariables}\}
        }{\indvariableof{\folexformula}} \, .
    \end{align*}
    We here denote the head variables of the basis encoding to $\bencodingof{\groundingof{\extformulaof{\atomenumerator}}}$ by $\catvariableof{\atomenumerator}$ to highlight their interpretation as propositional atoms.
\end{definition}

We depict the relation of a grounding tensor to a propositional formula as:
\begin{center}
    \input{tikz/syntax/propositionalization.tex}
\end{center}

\subsection{Quantifiers}

Existential and universal quantifiers appear in generic \firstOrderLogic{} and are besides substitutions further means to reduce the number of variables in a formula.
%They are not representable as linear transform of the respective quantifier-free formula.


% Definition of existential and universal quantifiction needed!
The semantics of existential quantification consists in a formula being true, if at least one state of the quantified variable is true, as we define next.

\begin{definition}
    Given a grounding tensor
    \begin{align*}
        \groundingofat{\folexformula}{\indvariableof{0},\ldots,\indvariableof{\indorder-1}} \,
    \end{align*}
    the existential and universal quantification with respect to the first variable are the tensors
    \begin{align*}
        \groundingofat{\left(\exists_{\indindexof{0}}\folexformula\right)}{\indvariableof{1},\ldots,\indvariableof{\indorder-1}} \andspace
        \groundingofat{\left(\forall_{\indindexof{0}}\folexformula\right)}{\indvariableof{1},\ldots,\indvariableof{\indorder-1}} \,
    \end{align*}
    with coordinates as follows.
    For an assignment $\indindexof{1},\ldots,\indindex$ to the non-quantified variables we have
    \begin{align*}
        \groundingofat{\left(\exists_{\indindexof{0}}\folexformula\right)}{\indexedindvariableof{1},\ldots,\indexedindvariableof{\indorder-1}} = 1
    \end{align*}
    if and only if there is an assignment $\indindexofin{0}$ such that
    \begin{align*}
        \groundingofat{\folexformula}{\indexedindvariableof{0},\indexedindvariableof{1},\ldots,\indexedindvariableof{\indorder-1}} = 1 \, .
    \end{align*}
    Conversely, we have for the universal quantification that
    \begin{align*}
        \groundingofat{\left(\forall_{\indindexof{0}}\folexformula\right)}{\indexedindvariableof{1},\ldots,\indexedindvariableof{\indorder-1}} = 1
    \end{align*}
    if and only if for any assignment $\indindexofin{0}$ we have
    \begin{align*}
        \groundingofat{\folexformula}{\indexedindvariableof{0},\indexedindvariableof{1},\ldots,\indexedindvariableof{\indorder-1}} = 1 \, .
    \end{align*}
\end{definition}


Let us now show, that existential and universal quantification are coordinatewise transforms (see \defref{def:coordinatewiseTransform}) of contracted grounding tensors.
To this end, let us introduce the greater-$z$ indicator $\greaterthanfunction{z}$, where $z\in\rr$, as the function
\begin{align*}
    \greaterthanfunction{z} : \rr \rightarrow \ozset
    \quad, \quad \greaterthanfunctionof{z}{x} =
    \begin{cases}
        1 & \ifspace x > z\\
        0 & \text{else}
    \end{cases} \, .
\end{align*}

\begin{theorem}
    For any formula $\folexformula$ with variables $\shortindvariablelist$ we have
    \begin{align*}
        \groundingofat{\left(\exists{\indindexof{0}}\folexformula\right)}{\indvariableof{1},\ldots,\indvariableof{\indorder-1}} =
        \coordinatetrafowrtofat{\existquanttrafo}{\contractionof{\groundingof{\folexformula}}{\indvariableof{1},\ldots,\indvariableof{\indorder-1}}}{\indvariableof{1},\ldots,\indvariableof{\indorder-1}}
    \end{align*}
    and
    \begin{align*}
        \groundingofat{\left(\forall{{\indindexof{0}}} \folexformula\right)}{\indvariableof{1},\ldots,\indvariableof{\indorder-1}}=
        \coordinatetrafowrtofat{\universalquanttrafo}{\contractionof{\groundingof{\folexformula}}{\indvariableof{1},\ldots,\indvariableof{\indorder-1}}}{\indvariableof{1},\ldots,\indvariableof{\indorder-1}}
    \end{align*}
\end{theorem}
\begin{proof}
    We proof the claimed equalities to arbitrary slices of the remaining variables, which amount to arbitrary substitutions of the formulas.
    For any indices $\indindexofin{1},\ldots,\indindexofin{\indorder-1}$ We notice that
    \begin{align*}
        \contractionof{\groundingof{\folexformula}}{\indexedindvariableof{1},\ldots,\indexedindvariableof{\indorder-1}}
        &= \sum_{\indindexofin{0}} \groundingofat{\folexformula}{\indexedindvariableof{0},\ldots,\indexedindvariableof{\indorder-1}} \\
        &= \cardof{\indindexofin{0} \wcols \groundingofat{\folexformula}{\indexedindvariableof{0},\ldots,\indexedindvariableof{\indorder-1}}=1} \, .
    \end{align*}
    We can thus understand the contracted grounding tensor as storing in its coordinates the count of the coordinate extensions to the zeroth variable, such that the grounding tensor is satisfied.
    This is analogous to our interpretation of contracted propositional formulas as world counts.
    From this it is obvious, that the existential quantification is satisfied, if the count is different from zero, which is captured by the coordinatewise transform with $\existquanttrafo$.
    We therefore arrive at
    \begin{align*}
        \groundingofat{\left(\exists_{\indindexof{0}}\folexformula\right)}{\indexedindvariableof{1},\ldots,\indexedindvariableof{\indorder-1}} =
        \coordinatetrafowrtofat{\existquanttrafo}{\contractionof{\groundingof{\folexformula}}{\indvariableof{1},\ldots,\indvariableof{\indorder-1}}}{\indexedindvariableof{1},\ldots,\indexedindvariableof{\indorder-1}} \, .
    \end{align*}
    The first claim follows, since the assignment to the non-quantified variables was arbitrary.
    The universal quantification is satisfied, when all extensions are satisfied, and the count is $\inddim$.
    Since $\inddim$ is the maximal count, this is captured by the coordinatewise transform with $\universalquanttrafo$ and we get
    \begin{align*}
        \groundingofat{\left(\forall{\indindexof{0}}\folexformula\right)}{\indexedindvariableof{1},\ldots,\indexedindvariableof{\indorder-1}} =
        \coordinatetrafowrtofat{\universalquanttrafo}{\contractionof{\groundingof{\folexformula}}{\indvariableof{1},\ldots,\indvariableof{\indorder-1}}}{\indexedindvariableof{1},\ldots,\indexedindvariableof{\indorder-1}} \, .
    \end{align*}
    With the same argument, the second claim is established.
\end{proof}

% Customized quantifiers
We can extend this discussion towards more generic counting quantifiers, of which the existential and the universal quantifier are extreme cases.
One can define quantifiers by demanding that at least $z\in\nn$ compatible groundings are satisfied, and show that they amount to coordinatewise transforms with $\greaterthanfunction{z}$.
What is more, quantifiers demanding that at most $z\in\nn$ are satisfied would be representable by transforms with an analogously defined function $\ones_{\leq z}$.
Such customized quantifiers appear for example in the $\mathrm{OWL\,2}$ standard of description logics (see \cite{rudolph_foundations_2011} and \secref{sec:kgRepresentation}).

% basis encodings
As will be discussed in \charef{cha:basisCalculus}, any coordinatewise transform can be performed by a contraction of a basis encoding of the tensor with a head vector prepared by the transform function (see \theref{the:tensorFunctionComposition}).
In the case here, a direct implementation would require a dimension of these head variables by $\inddim$, which can be infeasible when having large object sets.

% Prenex
To summarize, let us assume a formula is in its prenex normal form, that is a collection of quantifiers are acting on a quantifier free part.
We can represent its grounding tensor in three steps:
\begin{itemize}
    \item Instantiations of the atom groundings with the assigned variables, as contractions of the basis encoding of the world tensor with atom selecting tensors.
    \item Propositional formula acting on the head variables of the predicate instantiations, representing the connectives combining the formula.
    \item Quantifiers as a composition of contractions closing the quantified variable and coordinatewise transforms with the respective greater-than indicators.
\end{itemize}

\subsection{Storage in Basis CP Decomposition}\label{sec:basisCPgrounding}

In many situations, grounding cores are sparse and representations as single tensor cores comes with a drastic overhead.
We often encounter sparse grounding tensors, where the number of non-zero coordinates (to be investigated by basis CP ranks in \charef{cha:sparseRepresentation}) satisfies
\begin{align*}
    \sparsityof{\groundingof{\folexformula}} << \inddim^{\cardof{\indvariableof{\folexformula}}} \, .
\end{align*}
In this case, since most coordinates vanish, the basis CP decomposition (see \secref{sec:basisCP}) enables a representation of the grounding with significantly lower storage demand, see \theref{the:sparseBasisCP}.
This is particularly useful for representing large relational databases, where each object has only a few relations with others, while the majority of possible relations remains unsatisfied.
We depict such CP decomposition of a formula grounding in \theref{fig:groundingCP}.

% Standard KB Encoding and Assumptions
Most logical syntaxes exploit $\ell_0$-sparsity, explicitly storing only known assertions.
The interpretation of unspecified assertions depends on the underlying assumptions.
Under the Closed World Assumption, for example, all unspecified assertions are assumed to be false.

\begin{figure}[t]
    \begin{center}
        \input{tikz/syntax/grounding_decomposition.tex}
    \end{center}
    \caption{Basis CP Decomposition of the grounding of $\folexformula$, following the scheme of \theref{the:sparseBasisCP}.
    Instead of direct storage of the grounding tensor $\groundingof{\folexformula}$, the non-zero coordinates are enumerated by a variable $\datvariable$ and the corresponding coordinates stored in leg-matrices $\legcoreof{\folexformula,\indenumerator}$.}
    \label{fig:groundingCP}
\end{figure}

%\subsection{Summary}
%
%Let us summarize the tensor encodings of the representations of the different concepts, given a single \firstOrderLogic{} world:
%\begin{center}
%    \begin{tabular}{|p{\threecolumnwidth}|p{\threecolumnwidth}|p{\threecolumnwidth}|}
%        \hline
%        \textbf{Concept}        & \textbf{Representation}                       & \textbf{Decomposition}                             \\
%        \hline
%        Constant Symbol         & Basis vector                                  &                                                    \\
%        \hline
%        Function Symbol         & \BasisEncoding{}: Boolean and directed tensor &                                                    \\
%        Term                    & ""                                            & Acyclic tensor network of represented functions    \\
%        \hline
%        Predicate Symbol        & \CoordinateEncoding{}: Boolean tensor         &                                                    \\
%%        Relation & Boolean tensor \\
%        Quantifier-free Formula & ""                                            & Contraction of represented terms with relations    \\
%        Formula                 & ""                                            & Transform of corresponding quantifier-free formula \\
%        \hline
%    \end{tabular}
%\end{center}

\subsection{Example: Relational Databases}

A database is understood as a specific \firstOrderLogic{} world, and are operations on such a single world.
Queries are described by a formula $\impformula$, which are asked against a specific world $\worldindices$ to retrieve the grounding $\groundingof{\impformula}$.
The variables of such formulas are called projection variables.
The answer $\groundingof{\impformula}$ of a query is most conveniently represented as a list of solution mappings from the projection variables to objects in the world, such that the query formula is satisfied.
Answering a query by solution mappings corresponds with finding the basis CP Decomposition (see \secref{sec:basisCP}) of $\groundingof{\impformula}$.
We can understand these solution mappings as stored in the leg-matrices $\legcoreof{\folexformula,\indenumerator}$ (see \figref{fig:groundingCP}).

Let us give with the outer join an example of a popular operation to define queries, which efficient execution and storage can be improved based on considerations in the tensor network formalism.

\begin{definition}[Outer join]
    Let there be a world $\worldindices$ and formulas $\extformulaof{\selindex}$ depending on variables $\indvariableof{\nodesof{\selindex}}$, which have grounding tensors by
    \begin{align*}
        \groundingofat{\extformulaof{\selindex}}{\indvariableof{\node}} \wcols  \bigtimes_{\node\in\nodesof{\selindex}}[\inddimof{\node}] \rightarrow \ozset \, .
    \end{align*}
    Then their (outer) $\joinsymbol$ is defined as the grounding of their conjunctions, as
    \begin{align*}
        \groundingofat{\joinsymbol\left(\extformulaof{0},\ldots,\extformulaof{\seldim-1}\right)}{\bigcup_{\selindexin}\indvariableof{\nodesof{\selindex}}}
        = \contractionof{\groundingofat{\extformulaof{\selindex}}{\indvariableof{\nodesof{\selindex}}}\,:\,\selindexin}{\bigcup_{l\in[p]}\indvariableof{\nodesof{\selindex}}} \, .
    \end{align*}
\end{definition}

%Visualization and efficiency
We can understand the $\joinsymbol$ of groundings by a factor graph, where each grounding tensor decorates the hyperedge to the node set $\nodesof{\selindex}$.
The projection variable assignment to each formula combined in a $\joinsymbol$ operation provide a basic tensor network format to store the output of the operation.
There are thus situations, in which the solution map storage corresponding with a CP Decomposition comes with unnecessary overheads compared with other formats.

% Coordinatewise transform
We can also understand the $\joinsymbol$ operation as a coordinatewise transform (see \defref{def:coordinatewiseTransform}) with the product as transform function.
To make this connection solid, one would need to extend each joined formula trivially to the variables appearing in other formulas.

% Evaluation similar constraint propagation
The efficiency of evaluating the contraction to a $\joinsymbol$ operation might be improved by understanding it as an Constraint Satisfaction Problem (see \charef{cha:logicalReasoning}).
When applying efficient Message Passing algorithms such as Knowledge Propagation (see \algoref{alg:knowledgePropagation}), the groundings can be sparsified by local constraint propagation operations before turning to more global and more demanding contraction operations.
Here the groundings $\groundingof{\extformulaof{\selindex}}$ would be used to initialize Knowledge Cores $\kcoreof{\edge}$ and sequentially sparsified during the algorithm.

%\begin{example} % WOULD NEED OVERWORK: DRAW!
%	For example take a query with many basic graph patterns with pairwise different projection variables.
%	The global CP Decomposition would come here with an exponential storage overhead compared with storage as a tensor product of CP Decompositions to each Basic graph pattern.
%\end{example}

%% CONFUSING?
%\begin{remark}[Distinguishing from probabilistic queries]
%	Let us distinguish the discussion here from those of queries in probabilistic reasoning, which have two main differences.
%	First, we ask queries against all possible pairs of variables, instead of asking the probability of satisfaction of a specific formula.
%	Second, since we made the epistemologic assumption of knowing possibilities and not probabilities in logics, a query is answered by a truth value.
%	We then only output in the shape of solution mappings the variable assignments where the query formula is true.
% 	Thus, the queries here can be thought of as a batch of probabilistic queries with Boolean answers.
%	% Alternative -> Later?
%	Probabilistic queries can furthermore be understood in terms of the data extraction process described in this section.
%	We can ask the query in probabilistic form (decomposed into atomic formulas) on the resulting empirical distribution.
%	This results in the ratio of the worlds satisfying the query among those worlds satisfying the extraction query $\impformula$.
%\end{remark}




